\section*{Laws of Large Numbers}

\subsection*{11.1 Some Useful Bounds and Inequalities}

11Laws of Large Numbers

The laws of large numbers give operational meaning to the expectation of a random

variable. If we generate independent realizations of a random variable and compute

their mean, the weak law of large numbers states that the sample mean converges

to the expected value in probability as we increase the number of samples. In

this chapter, we present two versions of the weak law of large numbers. The first

one assumes that the random variables have finite variance, while the second one

assumes that the mean is finite.

The weak law has numerous applications. We discuss two sample applications:

Monte Carlo integration and data compression. Monte Carlo integration is a simple

but effective method for approximating integral in high dimensions, which arises

in financial computations. Data compression is a topic in information theory. The

fundamental result by Shannon on data compression is based on the weak law of

large numbers. We end this chapter with a version of the strong law of large number

that assumes the fourth moments of the random variables are finite.

The first inequality is the complex version of Cauchy-Schwarz inequality. In the

next theorem,X* denote the complex conjugate of X .

\begin{thmbox}{11.1 (Cauchy-Schwarz Inequality)}
\end{thmbox}

For complex-valued random variables X and Y in L2, the expectation

E[X*Y] is well-defined, and

\vertE[X*Y]\vert \leq

\sqrt

E[\vertX\vert2]E[\vertY\vert 2] (11.1)

holds, with equality holding iff X and Y are scalar multiple of each other with

probability 1. When X and Y are real, E[X*Y] is simplified to E[XY ].

\textit{Proof} Let X and Y be complex-valued random variables such that both E[\vertX\vert2]

and E[\vertY\vert2] are finite. We first show that \vertX*Y\vert is integrable. Let A1 be the subset

of \Omega in which \vertX()\vert > \vertY()\vert, and A2 be the subset in which \vertX()\vert\leq\vert Y()\vert.

Then, we have

E[\vertX*Y\vert] =

\int

A1

\vertX*Y\vert+

\int

A2

\vertX*Y\vert\leq

\int

A1

\vertX\vert2 +

\int

A2

\vertY\vert2 \leq E[\vertX\vert2]+ E[\vertY\vert2].

This shows that .

\int

\vertX*Y\vert is finite, and hence, X*Y is integrable.

If E[\vertX\vert2]= 0, then X = 0 almost surely, and hence, X*Y = 0 almost surely.

Both sides of (11.1) are equal to zero. Suppose E[\vertX\vert2] > 0. Define a function

f( z )\coloneqqE[\vertzX + Y\vert2] with z being a complex variable. We can expand it as

f( z )= E[(zX +Y)*(zX +Y)]=\vert z\vert2E[\vertX\vert2]+ z*E[X*Y]+ zE[XY *]+ E[\vertY\vert2].

By completing square, we obtain

f( z )= E[\vertX\vert2]\cdot

\vert\vert\vertz+ E[X*Y]

E[\vertX\vert2]

\vert\vert\vert

- \vertE[X*Y]\vert2

E[\vertX\vert2] +E [\vertY\vert 2].

Since f( z )\geq 0 for all z, the constant term on the right must be nonnegative. This

yields the Cauchy-Schwarz inequality.

When equality holds, then f( z0) = 0 for a particular choice of z0. This implies

thatE[\vertz0X + Y\vert2]= 0, and hence, Y =- z0X almost surely. \hfill $\square$

Note that the proof of Theorem 11.1 does not rely on the assumption that the

underlying measure space is a probability space. The Cauchy-Schwarz inequality

can be derived for a general measure, as long as X and Y are square integrable.

\begin{thmbox}{11.2 (Chebyshev Inequality)}
\end{thmbox}

If the variance of a real-valued random variable X is finite, then for any

\epsilon> 0,

P(\vertX - E[X]\vert \geq \epsilon)\leq Va r(X)

\epsilon2 .

\textit{Proof} The indicator function .1{\vertX-E[X]\vert\geq\epsilon}() satisfies

1{\vertX-E[X]\vert\geq\epsilon}()\leq (X() - E [X])2/\epsilon2 for all \in \Omega.

Taking the expectation of both sides, we obtain

P( \vertX- E [X]\vert \geq\epsilon)= E [1{\vertX-E[X]\vert\geq\epsilon}] \leq E[(X- E [X])2]

\epsilon2 = Va r(X)

\epsilon2 .

\hfill $\square$

The next result is called Jensen inequality .

\begin{thmbox}{11.3 (Jensen Inequality)}
\end{thmbox}

Suppose \phi is a convex function and X is a real-valued random variable with

finite mean. Then E[\phi(X)]\geq \phi(E[X]).

\textit{Proof} By the convexity of \phi, at each point x* in the domain of \phi, we can draw a

straight line that touches the graph of \phi at one point, ie., it passes through the point

(x*,\phi( x*)) and lies beneath the graph of \phi(x).

We pick x* =E [X], which is finite by assumption. Let g(x)= ax + b be a

linear function such that g(E[X])= \phi(E [X]) andg(x)\leq \phi(x) for all x I f \phi is not

differentiable at x* =E [X], there could be more than one choice of coefficients a

and b In this case, we just pick one arbitrarily.

By applying the monotonic property of integral, we obtain

E[\phi(X)]\geq E[g(X)]= aE[X]+ bE[1]= aE[X]+ b= \phi(E [X]).

Note that we have used E[1]= 1 in the second last equality. \hfill $\square$
